% Standalone audit. Compile this file directly with pdflatex (two passes).
% No existing project file, proof, status, or compendium is changed by A01.
% CORE VERDICT: VERIFIED
% EXTENDED VERDICT: MINOR REVISION
\documentclass[11pt]{article}
\usepackage[T1]{fontenc}
\usepackage[utf8]{inputenc}
\usepackage{lmodern,amsmath,amssymb,amsthm}
\usepackage[margin=0.9in]{geometry}
\usepackage{array,booktabs,longtable,enumitem}
\usepackage[hidelinks]{hyperref}
\newcolumntype{P}[1]{>{\raggedright\arraybackslash}p{#1}}
\setlength{\emergencystretch}{3em}
\setlength{\parindent}{0pt}
\setlength{\parskip}{5pt}
\setlist{nosep,leftmargin=*}
\newcommand{\RR}{\mathbb R}
\newcommand{\ZZ}{\mathbb Z}
\newcommand{\Orth}{\RR_{\geq0}^{n}}
\newcommand{\Dom}{(-1,\infty)^n}
\newcommand{\dist}{\operatorname{dist}_{\infty}}
\newcommand{\pos}[1]{(#1)_+}
\newcommand{\lin}[1]{\textbf{L#1}}
\newcommand{\pass}{\textbf{PASS}}
\newcommand{\minor}{\textbf{MINOR}}
\title{A01: Independent audit of P01\\Well-posedness and the absorbing box}
\author{Independent mathematical reconstruction and adversarial review}
\date{27 September 2026}
\begin{document}
\maketitle

\section{Verdicts, exact version, and scope}

\textbf{CORE VERDICT: VERIFIED.}
All nine requested canonical claims hold for the demand-modulated equation
with $n\geq5$, $1\leq q\leq n-1$, and $a,k,c,g,h>0$. In particular,
the smaller radius in P01 is correct, and the closed cube at that radius
is entered in finite time with one bound for every finite initial state
in the entire unbounded nonnegative orthant.

\textbf{EXTENDED VERDICT: MINOR REVISION.}
The weaker-parameter existence, positivity, comparison, zero-radius,
$k=0$, and exact absorption-classification results all pass. The sole
required correction is a root-description error in \lin{428}: for
$c=0$ and $a>0$, the archived quadratic has two nonnegative roots, so
$B$ is not unambiguously described as ``the nonnegative root.'' Replace
that phrase with ``the larger real root, which is nonnegative.'' The
displayed formula for $B$ and every conclusion using it are correct.
This is a minor defect of the exact extended text, not a failure of its
absorption theorem. No mathematical counterexample to a stated theorem
was found, and the verdicts rest on proofs below, not that observation.

\textbf{Audited source:}
\path{proofs/P01_wellposedness_and_absorbing_box.tex}, exactly 512 lines
in the current \path{cyclic_network_project} directory. Every reference
of the form L$x$--$y$ below refers to this source, including blank and
comment lines. The source SHA-256 is
\begin{quote}\small\ttfamily
ba881a703b7c1dedefce4ccb920f7b8243\allowbreak
418cccd9219eb8285fbb707911941b
\end{quote}
The source is byte-identical to its member in
\path{cyclic_network_project_P01_draft.zip}. The attached and local
\path{proofs_compendium.pdf} are byte-identical, with SHA-256
\begin{quote}\small\ttfamily
d696801db40fddc12aae7c2b8de43361ee0\allowbreak
991c6de5d186f6d1eada54e8faf61
\end{quote}
P01 occupies PDF pages 11--17. Its exact absorption result is
\emph{Theorem 5.6}; the later subsection also numbered 5.6 is the
checks subsection and must not be confused with that theorem.

The mathematical starting point is the canonical differential equation,
not an archived derivation or an internal validation report. The source
was checked by independently reconstructing each argument. Only ordinary
local ODE existence, uniqueness, continuous dependence, and continuation
for a locally Lipschitz vector field are imported. Finite-maximum,
comparison, positivity, and absorption arguments are supplied below.
No literature survey, novelty assessment, proof of P02, or analysis of
any later module is part of this audit. Dependency comments are conditional
consequences only. Existing files and their DRAFT/TODO status are unchanged.

\section{Equation and quantified assumptions}
The only vector field audited is
\begin{equation}\label{A01:field}
 F_i(x)=a x_i-kx_i^2+\frac{c x_{i-1}}{1+x_i}
       -\frac{g x_i}{1+x_{i+q}}-h x_i x_{i-q},
 \qquad i\in\ZZ/n\ZZ.
\end{equation}
All occurrences remain present when offsets coincide. The open local
domain is $\Dom$; forward boundedness claims concern $\Orth$ only.
For the extended results, write
\begin{equation}\label{A01:weak}
 n\in\ZZ_{\geq1},\qquad q\in\ZZ/n\ZZ,\qquad
 a\in\RR,\quad k>0,\quad c,g,h\geq0.
\end{equation}
Parameters and dimension are fixed and finite. Uniformity in initial
states does not mean uniformity in these parameters or in $n$.
Set $X=\max_i x_i$, $K_r=[0,r]^n$, $b(s)=a-ks+c/(1+s)$ and
$H(s)=s b(s)$. For nonnegative states,
\begin{equation}\label{A01:distance}
 \dist(x,K_r)=\pos{X-r}.
\end{equation}
Indeed, clipping each coordinate at $r$ attains this distance, and a
coordinate attaining $X$ gives the matching lower bound.

\section{Local solutions, faces, and positivity}
\subsection{Finite maxima and local well-posedness}
For a finite family of $C^1$ functions $y_j$ and $v=\max_j y_j$, all
inactive indices at a fixed time have a positive gap. They remain
inactive for a sufficiently short right interval. For the active set
$I(t)$, the finitely many expansions
$y_j(t+\varepsilon)=v(t)+\varepsilon\dot y_j(t)+o(\varepsilon)$
have a common $o(\varepsilon)$ error bound. Therefore
\begin{equation}\label{A01:max}
 \frac{d^+v}{dt}=D^+v(t)=\max_{j\in I(t)}\dot y_j(t).
\end{equation}
The right derivative exists; a two-sided derivative need not. Bounded
derivatives on a compact time interval make $v$ Lipschitz there.
This proves P01's lemma at unique maxima, ties, and switching times.
It makes no assertion that a single maximizing index can be continued
differentiably through a switch.

Every denominator in \eqref{A01:field} is positive on $\Dom$, including
when several denominators or coordinates coincide. Thus $F$ is smooth
there for arbitrary finite real coefficients. The standard local ODE
theorem gives a unique maximal classical solution and continuous
dependence on initial data on compact common existence intervals.
Smoothness on this open domain alone is not a global-existence assertion.

\subsection{Nonnegative invariance, including tangent faces}
If $x\in\Orth$ and $x_i=0$, then $F_i(x)=c x_{i-1}\geq0$ for $c\geq0$.
This sign alone needs a barrier argument at tangent faces. On a compact
interval inside the maximal existence interval, let $x^+$ be the
coordinatewise positive part and $V=\max(0,-x_0,\ldots,-x_{n-1})$.
The curves and their connecting segments lie in a compact rectangle
inside $\Dom$, where $F$ has a maximum-norm Lipschitz constant $L$.
If $V>0$ and $-x_i=V$, then $x_i^+=0$ and
\[
 -F_i(x)\leq |F_i(x)-F_i(x^+)|\leq LV.
\]
If $V=0$, active negative-coordinate derivatives are nonpositive and
the constant zero has derivative zero. Equation \eqref{A01:max} gives
$D^+V\leq LV$ in both cases. Absolute continuity and $V(0)=0$ give
$V\equiv0$. This does not assume global existence and does not impose
signs on $a,k,g,h$. The argument is valid on every compact subinterval
of the maximal interval.

\subsection{Integrating factors and arbitrarily early propagation}
Along the now nonnegative solution, define
\[
 B_i(t)=a-kx_i(t)-\frac{g}{1+x_{i+q}(t)}-h x_{i-q}(t),\qquad
 D_i(t)=\frac{c x_{i-1}(t)}{1+x_i(t)}.
\]
Then $\dot x_i=B_i(t)x_i+D_i(t)$ exactly, $D_i\geq0$, and the coefficients
are continuous on each compact existence interval. Their dependence on
the solution does not invalidate variation of constants:
\begin{equation}\label{A01:positive}
 x_i(t)=x_i(0)e^{\int_0^t B_i(v)\,dv}
       +\int_0^t e^{\int_s^t B_i(v)\,dv}D_i(s)\,ds.
\end{equation}
An initially positive coordinate stays positive. If $c>0$ and coordinate
$j$ is positive for every $s>0$, then $D_{j+1}(s)>0$ on every interval
$(0,t)$, so $x_{j+1}(t)>0$ for every $t>0$. Induct on the number of
step-one edges, proving a statement for \emph{all} positive times at
each induction step. This incurs no positive waiting time per edge.
At most $n-1$ steps reach every node. For $n=1$, no propagation is
needed; a nonzero nonnegative state is already positive. Neither $q$
nor $\gcd(n,q)$ appears in the propagation path.

For $c=0$, the integral vanishes and every zero coordinate remains zero.
Also $F(0)=0$, so zero initial data stay zero. These conclusions preclude
finite-time arrival at the origin from any nonzero nonnegative state.
They do not imply a positive lower bound as $t\to\infty$.

\section{Scalar comparison, radius, and quantitative attraction}
\subsection{The sign of each loss and the root algebra}
At every active maximum $x_i=X\geq0$,
\[
 \frac{c x_{i-1}}{1+X}\leq\frac{cX}{1+X},\qquad
 -\frac{gX}{1+x_{i+q}}\leq-\frac{gX}{1+X},\qquad
 -hXx_{i-q}\leq0.
\]
In the middle inequality, $1+x_{i+q}\leq1+X$ first makes the reciprocal
larger; multiplication by $-gX\leq0$ reverses the order. Thus
\begin{equation}\label{A01:dini}
 D^+X=\max_{i:x_i=X}F_i(x)
       \leq H(X)-\frac{gX}{1+X}\leq H(X).
\end{equation}
All active indices satisfy the same bound. No independent choice of
neighbors is used, so all coincident-offset cases are included.

Since $b'(s)=-k-c/(1+s)^2<0$, $b(0)=a+c$, and $b(s)\to-\infty$,
there is one positive root precisely when $a+c>0$. Multiplying its
equation by $1+R$ gives
\[
 (a-kR)(1+R)+c=0,
 \qquad kR^2+(k-a)R-(a+c)=0.
\]
The discriminant is
\[
 (k-a)^2+4k(a+c)=(a+k)^2+4kc.
\]
For $a+c>0$, the constant coefficient is negative, so the positive
root is the plus-sign root. For $a+c\leq0$, $k-a>0$ and
$(a+k)^2+4kc\leq(k-a)^2$, making that raw root nonpositive. Hence
\begin{equation}\label{A01:R}
 R=\max\left\{0,\frac{a-k+\sqrt{(a+k)^2+4kc}}{2k}\right\}.
\end{equation}
For $R>0$, $b$ is positive on $[0,R)$ and negative on $(R,\infty)$.
An exact identity, stronger than the stated inequality, is
\begin{equation}\label{A01:factor}
 H(s)=-s(s-R)\left[k+\frac{c}{(1+s)(1+R)}\right]
       \leq-ks(s-R)\quad(s\geq R).
\end{equation}
For $R=0$, $b(s)\leq a+c-ks\leq-ks$, so $H(s)\leq-ks^2$.
The selected zero radius need not satisfy the quadratic when $a+c<0$;
it is nevertheless a zero of $H$ because of the factor $s$.

\subsection{Scalar domination without a cooperativity assumption}
Let $u'=H(u)$, $u(0)=X(0)$. The scalar interval
$[0,\max\{X(0),R\}]$ is invariant by local uniqueness and the endpoint
signs. Consequently this scalar solution is global. On any compact
interval before the vector solution's maximal time, put
$W=(X-u)_+=\max(0,x_i-u:i\in\ZZ/n\ZZ)$. If $W>0$,
\[
 D^+W\leq H(X)-H(u)\leq L(X-u)=LW
\]
for a scalar Lipschitz constant $L$. At $W=0$, an active coordinate
requires $x_i=X=u$ and gives $D^+(x_i-u)\leq0$; if there is none, the
constant zero is the only active function. Thus $D^+W\leq LW$
everywhere. Absolute continuity gives $W=0$. In particular,
\begin{equation}\label{A01:bound}
 0\leq X(t)\leq u(t)\leq\max\{X(0),R\}.
\end{equation}
No monotonicity of $H$ and no order-preserving property of the full
network are needed. The inhibitory couplings therefore cause no gap
in this scalar comparison.

If the vector maximal time were finite, \eqref{A01:bound} would confine
the solution to a compact cube in $\Dom$, separated from every singular
hyperplane by at least one. The vector field is bounded on a compact
neighborhood of that cube, so $x$ has a limit there at the maximal time.
Local continuation at that limit is a contradiction. This proves global
existence and also proves positive invariance of every $K_r$, $r\geq R$.

\subsection{Rates, uniformity over unbounded data, and outer cubes}
For $R>0$, \eqref{A01:factor} yields
\begin{equation}\label{A01:rate}
 \dist(x(t),K_R)\leq e^{-kRt}\pos{X(0)-R}.
\end{equation}
This estimate alone gives uniform attraction only on bounded initial
sets. The stronger assertion requires a second calculation. If $u>R$,
put $w=1-R/u$. Then
\[
 w'=\frac{Ru'}{u^2}\leq-kR w,
 \qquad 0<w(0)<1.
\]
Thus, for $u(0)>R$,
\[
 u(t)\leq\frac{R}{1-(1-R/u(0))e^{-kRt}}
          \leq\frac{R}{1-e^{-kRt}}\quad(t>0).
\]
The last bound also holds for initial data at or below $R$. Together
with \eqref{A01:distance}, it proves the exact strong quantifier
\begin{equation}\label{A01:uniform}
 \sup_{x^0\in\Orth}\dist(x(t;x^0),K_R)
 \leq\frac{R}{e^{kRt}-1}\longrightarrow0\quad(t\to\infty).
\end{equation}
For $R=0$, integration of $(1/u)'\geq k$ wherever $u>0$ gives
\begin{equation}\label{A01:zero}
 X(t)\leq\frac{X(0)}{1+kX(0)t},\qquad
 \sup_{x^0\in\Orth}X(t;x^0)\leq\frac1{kt}\quad(t>0).
\end{equation}
Zero initial data are handled separately by uniqueness. These suprema
are over finite vectors of arbitrary size, not an initial state at
infinity. Their bounds need not be finite at $t=0$.

Solving $R/(1-e^{-kRt})\leq r$ for $r>R$ gives
\[
 T_r=\frac1{kR}\log\frac{r}{r-R}\quad(R>0),\qquad
 T_r=\frac1{kr}\quad(R=0).
\]
Every solution is in $K_r$ at this common time, and positive invariance
keeps it there. All constants and inequality directions in the scalar
comparison portion of P01 pass.

\section{Exact finite-time absorption: both directions}
Assume \eqref{A01:weak} and first take $R>0$. Set
\[
 t_0=\frac{\log2}{kR},\qquad M=2R.
\]
The preceding initial-state-independent bound puts every trajectory
in the invariant cube $K_M$ by $t_0$.

\subsection{The direct strict-loss cases}
If $g>0$, then at all times when $X\geq R$,
\[
 D^+X\leq-\frac{gX}{1+X}
          \leq-\frac{gR}{1+R}=-\delta,
 \qquad\delta=\frac{gR}{1+R}>0.
\]
Integrate the almost-everywhere derivative inequality until first entry.
For $X(0)\leq M_0$, entry occurs by
\[
 \frac{\pos{M_0-R}}{\delta}
   =\frac{\pos{M_0-R}(1+R)}{gR}.
\]
Applying this after $t_0$ with $M_0=2R$ proves
\begin{equation}\label{A01:Tstar}
 T_*=\frac{\log2}{kR}+\frac{1+R}{g}.
\end{equation}
If entry had not occurred at the endpoint of either time bound,
integration and continuity would give $X\leq R$ there, a contradiction.
Thus ``by'' includes equality of times. This supplies the canonical
answer, without using $h>0$ or positivity propagation.

If $g=0$, $h>0$, and $q\equiv0\pmod n$, then $x_{i-q}=X$ at a maximum.
Consequently $D^+X\leq H(X)-hX^2\leq-hR^2$ outside $K_R$.
After $t_0$, at most another $R/(hR^2)=1/(hR)$ is required.

\subsection{The propagated strict-loss case and every constant}
Suppose $g=0$, $h>0$, $c>0$; retaining $g$ in the next constant is
harmless and also makes the lower inequality valid for $g\geq0$.
Choose any one fixed $\tau>0$ and define
\begin{equation}\label{A01:constants}
 \beta=(-a)_++(k+h)M+g,\qquad
 \eta=\frac{c}{1+M}>0,\qquad M=2R.
\end{equation}
On $K_M$, each negative coefficient is bounded in the required lower
direction: $a\geq-(-a)_+$, $-kx_i\geq-kM$,
$-g/(1+x_{i+q})\geq-g$, and $-hx_{i-q}\geq-hM$. Also
$c/(1+x_i)\geq c/(1+M)$. Multiplying these estimates by nonnegative
coordinates proves
\begin{equation}\label{A01:lower}
 \dot x_i\geq-\beta x_i+\eta x_{i-1}.
\end{equation}
No omitted dependence on $q$ is needed in $\beta$ or $\eta$.

For $s\geq t_0$, set $z_i(r)=e^{\beta(r-s)}x_i(r)$. Then
$z_i'\geq\eta z_{i-1}$ and
$z_i(r)\geq x_i(s)+\eta\int_s^r z_{i-1}(v)\,dv$.
For any integer $m\geq0$, iterative substitution gives
\[
 z_i(s+\tau)\geq
 \sum_{j=0}^{m}\frac{(\eta\tau)^j}{j!}x_{i-j}(s)
       +\mathcal E_{m+1},\qquad \mathcal E_{m+1}\geq0.
\]
Here $\mathcal E_{m+1}$ is the $(m+1)$-fold time-ordered integral of
$z_{i-m-1}$ with coefficient $\eta^{m+1}$. Its sign is nonnegative
even after the indices wrap around the cycle. Taking $m=n-1$ proves
\begin{equation}\label{A01:series}
 x_i(s+\tau)\geq e^{-\beta\tau}
    \sum_{j=0}^{n-1}\frac{(\eta\tau)^j}{j!}x_{i-j}(s).
\end{equation}
The highest retained power is $n-1$, not $n$; the unused remainder
has $n$ integrals. This resolves any ambiguity in ``$n$-fold
propagation.'' For $n=1$, the retained sum is just $x_i(s)$.

If a trajectory has not entered $K_R$ by $t\geq t_0+\tau$, then
$X(t-\tau)>R$. For each target $i$, the list $i-j$, $0\leq j<n$,
contains an index achieving that earlier maximum. One term of
\eqref{A01:series} therefore gives
\begin{equation}\label{A01:d}
 x_i(t)\geq d:=R e^{-\beta\tau}
       \min_{0\leq j<n}\frac{(\eta\tau)^j}{j!}>0.
\end{equation}
This is uniform in the initial state, conditional on nonentry. It is
\emph{not} an unconditional lower bound on all of $\Orth$: the zero
solution already disproves such a claim. The proof needs only this
conditional floor, at each time up to first entry. Since $X\geq R$
and every $x_{i-q}\geq d$ at those times,
\[
 D^+X\leq H(X)-hXx_{i-q}\leq-hRd.
\]
At $t_0+\tau$ the maximum is at most $2R$; thus the remaining excess
is at most $R$ and the additional time is at most
$R/(hRd)=1/(hd)$. The common bound is exactly
\begin{equation}\label{A01:Th}
 T_h=t_0+\tau+\frac1{hd}.
\end{equation}
All of $\beta,\eta,d,\tau$ and this final cancellation are correct.
Every fixed $\tau>0$ works. For example, $\tau=1/\beta$ is permissible
because $\beta>0$ when $R>0$ and $k>0$, but no optimization is needed.

\subsection{Necessity and logical equivalence}
Under nonnegative $c,g,h$, failure of
\begin{equation}\label{A01:criterion}
 g>0\quad\text{or}\quad
 [h>0\text{ and }(c>0\text{ or }q\equiv0\pmod n)]
\end{equation}
has exactly two cases:
\begin{enumerate}
\item $g=h=0$. Direct substitution of $x_i=s$ in the original equation
gives $s'=H(s)$. For $s(0)>R$, the solution cannot meet the equilibrium
$R$ in finite time by local uniqueness, also applied backward near a
putative meeting. It stays above $R$ forever at finite times.
\item $g=0$, $h>0$, $c=0$, and $q\not\equiv0\pmod n$.
Then $R=a/k>0$. Give one coordinate $j$ initial value $s_0>R$ and
all others value zero. Zero coordinates stay zero, including $j-q$.
The surviving coordinate obeys $s'=as-ks^2$ and
\[
 s(t)=\frac{R}{1-(1-R/s_0)e^{-at}}>R\quad(t<\infty).
\]
This is a direct counterexample to absorption, for every such $n,q$.
\end{enumerate}
Thus failure of \eqref{A01:criterion} prevents even pointwise absorption
of all trajectories. Sufficiency above gives one time for the entire
orthant. That implies absorption of bounded subsets, which implies
pointwise absorption by considering singleton sets. These implications,
combined with the two counterexamples, prove the equivalence of all four
items in P01's Theorem 5.6. Pointwise absorption is not being upgraded
to uniform absorption by a general compactness principle.

When $R=0$, \eqref{A01:positive} prevents any nonzero solution from
entering $\{0\}$ at a finite time. The attraction estimates still hold.
Even if \eqref{A01:criterion} is true, it cannot override the explicit
$R>0$ hypothesis in the classification.

\section{Weaker parameters, degeneracies, and signed counterexamples}

\subsection{Small dimensions and coincident offsets}
For $n=1$, every residue is zero and the original equation is exactly
\[
 \dot x=x\left[a-kx+\frac{c-g}{1+x}-hx\right].
\]
The absorption criterion for $R>0$ reduces to $g>0$ or $h>0$.
If either is positive, the scalar field at $R$ has a strict negative
gap. If both vanish, it is exactly $H$ and a solution above $R$ never
meets $R$ in finite time. Thus the classification remains correct;
the sparse nonzero-offset counterexample is simply inapplicable here.

For $n=2$, $q=1$, predecessor, successor, and both step-$q$ neighbors
coincide. Every original term is still present and all three inequalities
leading to \eqref{A01:dini} remain true simultaneously. For even $n$
and $q=n/2$, the two inhibitory offsets coincide. Again their contributions
add rather than disappear. When $q=1$ or $q=n-1$, one inhibitory neighbor
coincides with the predecessor. No estimate assumed distinct coordinates.
For $n=6$, $q=2$, the step-$q$ graph is disconnected, but the positive
input still uses the connected step-one cycle. The same applies to
every noncoprime pair $(n,q)$.

\subsection{Zero and small coefficients}
If $c=0$, no new positive coordinate can be created. For $R>0$,
$R=a/k$, and absorption occurs exactly when $g>0$ or when both $h>0$
and $q\equiv0$. If $g=0$ but $h,c>0$, the propagation proof gives
absorption for every residue. If $h=0$, it occurs exactly when $g>0$.
If $g=h=0$, a synchronous scalar trajectory supplies nonabsorption,
even for $c>0$ and strictly positive data.

At $a+c=0$, direct substitution gives
\[
 H(s)=-ks^2-\frac{cs^2}{1+s}\leq-ks^2,
\]
so there is no missing equality case. For $a+c<0$ the extra linear
decay is favorable. In both cases $R=0$ and there is no finite-time
entry from nonzero data. All expressions containing $1/R$ in P01 are
explicitly restricted to $R>0$.

For arbitrarily small positive $R$, all fixed-parameter proofs still
hold. One exact test family is to prescribe $R=\varepsilon>0$ and set
$a=k\varepsilon-c/(1+\varepsilon)$. Then $a+c>0$ and $b(\varepsilon)=0$
exactly. To keep the whole tuple canonical, one may choose
$k=1$, $a=c=\varepsilon>0$, with $g,h>0$; its positive radius tends
to zero as $\varepsilon\downarrow0$. The displayed entry times need
not remain bounded in that limit. Likewise the propagated floor may
tend to zero as $c\downarrow0$, and the resulting bound may diverge
as $h\downarrow0$ or $n$ increases. These are parameter limits, not
counterexamples to uniformity in initial data at fixed parameters.

For numerically evaluating a small positive root when $a<k$, a
cancellation-resistant equivalent formula is
\[
 R=\frac{2(a+c)}{\sqrt{(a+k)^2+4kc}+k-a}\qquad(a+c>0).
\]
This is an optional numerical presentation improvement, not a correction
of P01's exact root formula. No numerical root computation is used to
justify the theorem.

\subsection{The \texorpdfstring{$k=0$}{k=0} extension}
For $k\geq0$, $a\in\RR$, $c,g,h\geq0$, dropping nonpositive terms at
a maximum gives
\[
 D^+X\leq aX+\frac{cX}{1+X}\leq(a_++c)X.
\]
The positive-part comparison argument therefore proves
$X(t)\leq X(0)e^{(a_++c)t}$. This is enough to exclude escape from
compact subsets of $\Dom$ before any finite maximal time, proving
global existence. The previous face and integrating-factor arguments
remain valid. Strict propagation still requires $c>0$ (unless $n=1$).

No claim about $R$ or an absorbing cube is imported into $k=0$.
With $a>0$ and $k=c=g=h=0$, $x_i(t)=x_i(0)e^{at}$ is unbounded for
any positive coordinate. This refutes a bounded-absorption guarantee
over that class, without asserting that every $k=0$ system is unbounded.

\subsection{All stated signed-parameter failures}
\begin{enumerate}
\item \textbf{Negative $c$.} If $n\geq2$, select $x_i=0$ and
$x_{i-1}=1$, with other coordinates nonnegative. Then $F_i=c<0$ and
$x_i(t)=ct+o(t)<0$ for small $t>0$. The restriction $n\geq2$ matters:
for $n=1$, $x_i=0$ also forces its predecessor to be zero, and every
term in the scalar equation has a factor $x$. The stated example does
not claim a failure for $n=1$.
\item \textbf{Negative $k$.} With $a=c=g=h=0$ and one coordinate
$s_0>0$, its exact solution is
$s(t)=s_0/(1-|k|s_0t)$, which blows up at $1/(|k|s_0)$.
This works for every dimension and residue.
\item \textbf{Negative $h$.} With $k>0$, $a=c=g=0$, $h<-k$,
and all initial coordinates $s_0>0$, direct substitution gives
$s'=-(k+h)s^2$. Thus
$s(t)=s_0/[1-(-(k+h))s_0t]$ blows up in finite time. Coincident
offsets do not alter this calculation.
\end{enumerate}
All three examples are valid existential failures outside the proved
parameter domain. They do not assert necessity of every sign assumption
for every individual signed coefficient tuple.

\subsection{Archived radius: the single required correction}
The archived formula
\[
 B=\frac{a+\sqrt{a^2+4kc}}{2k}
\]
is always the larger real root of $as-ks^2+c=0$. When $c>0$, $B>0$
and the root relation yields
\[
 H(B)=(aB-kB^2)+\frac{cB}{1+B}
       =-c+\frac{cB}{1+B}=-\frac{c}{1+B}<0.
\]
If $R>0$, this and the strict decrease of $b$ imply $B>R$; if $R=0$,
it follows directly from $B>0$. Thus its cube is a strictly larger
absorbing cube, even when the exact-radius cube is nonabsorbing.
For $c=0$, $B=\max(a,0)/k=R$. All these statements pass.

\textbf{Issue E01 (minor; \lin{428}).} The wording ``the nonnegative
root'' implicitly identifies a unique such root. For $a=k=1$ and
$c=0$, the polynomial is $s-s^2$ and its two nonnegative roots are
$0$ and $1$, whereas the displayed $B$ selects $1$. Suggested replacement:
\begin{quote}
``is the larger real root of $as-ks^2+c=0$, and is nonnegative.''
\end{quote}
This is the only required correction found. No formula, inequality,
theorem hypothesis, time constant, or theorem conclusion needs changing.
The CORE verdict is unaffected because its $c>0$ case has a unique
nonnegative root. The EXTENDED verdict is conservatively held at
MINOR REVISION until the exact text is corrected; no P01 edit is made here.

\section{Terminology and quantifiers}
Let $\Phi_t(x^0)$ denote the established nonnegative forward solution.
The following distinctions are essential; equivalences here are specific
to this equation and parameter class, not general facts about flows.

\begin{longtable}{@{}P{0.24\textwidth}P{0.71\textwidth}@{}}
\toprule
Property & Exact meaning and audited conclusion \\
\midrule\endhead
Positive invariance & $\Phi_t(K_R)\subseteq K_R$ for every $t\geq0$.
Holds under \eqref{A01:weak}, including $R=0$. This does not assert
$\Phi_t(K_R)=K_R$ or backward invariance. \\
Pointwise absorption & For each $x^0$, some finite $T(x^0)$ places the
solution in $K_R$ for all later times. For $R>0$, holds exactly under
\eqref{A01:criterion}. For $R=0$, fails for every nonzero initial state. \\
Absorption of bounded sets & For each bounded $E\subset\Orth$, one
$T_E$ works for all its states. Equivalent to the preceding property
in the exact classification, because sufficiency provides a stronger bound. \\
One time for the orthant & There exists $T<\infty$ such that
$\Phi_t(\Orth)\subseteq K_R$ for all $t\geq T$. Holds under the criterion
when $R>0$; $T_*$ works canonically. Parameters are fixed. \\
Global attraction & For every initial state, the distance to $K_R$
tends to zero. Holds under \eqref{A01:weak}, irrespective of the exact
absorption criterion. \\
Uniform attraction on bounded sets &
$\sup_{x^0\in E}\dist(\Phi_t(x^0),K_R)\to0$ for each bounded $E$.
Follows from \eqref{A01:rate} or \eqref{A01:zero}. \\
Uniform attraction on the orthant &
$\sup_{x^0\in\Orth}\dist(\Phi_t(x^0),K_R)\to0$.
Follows from \eqref{A01:uniform} and \eqref{A01:zero}, whose right-hand
sides contain no $X(0)$. This is valid even if exact-radius absorption fails. \\
Lyapunov stability as a set & Relative to $\Orth$, for every
$\varepsilon>0$, taking $\delta=\varepsilon$ suffices: initial distance
less than $\delta$ remains less than $\varepsilon$. The distance bounds
never exceed the initial distance. The statement concerns distance to
the whole cube, not motion within it. \\
Global asymptotic stability as a set & Under P01's explicit convention
of Lyapunov stability plus attraction, holds relative to $\Orth$.
It does not imply convergence to one equilibrium, asymptotic stability
of each state in the cube, or minimality as a global attractor. \\
\bottomrule
\end{longtable}

If a different convention requires strict two-sided invariance for an
``asymptotically stable set,'' retain the two explicit properties instead
of changing the mathematics. P01 already states its convention, so this
is not an additional defect. The direction of the set distance is also
important: attraction here is the supremum of point-to-cube distances,
not convergence in symmetric Hausdorff distance of images to the entire cube.

\section{Line-by-line issue and verification register}
The following consecutive source ranges cover all 512 lines, grouping
only adjacent lines belonging to one proof step. PASS records a checked
step; it is not inferred from compilation. E01 is the sole required
revision. Explanatory cautions already satisfied by the text are not
reported as missing hypotheses or errors.

\small
\begin{longtable}{@{}P{0.12\textwidth}P{0.12\textwidth}P{0.71\textwidth}@{}}
\toprule
P01 lines & Finding & Independent check / disposition \\
\midrule\endhead
1--8 & \pass & Metadata correctly confines P01 to the canonical equation
and standard ODE inputs. DRAFT status is not evidence of correctness and
is not changed by this audit. \\
9--24 & \pass & The weaker scope admits every residue, including zero,
and nonnegative $c,g,h$. Indices presuppose an integer $n\geq1$.
All coincident terms must remain. \\
25--38 & \pass & Maximum, positive part, Dini derivative, and maximum-norm
distance formula are correct; the last uses $x\in\Orth$. \\
39--47 & \pass & Bounded-set absorption and uniform attraction have the
right quantifiers. The stronger whole-orthant assertion is separately proved. \\
48--61 & \pass & A finite maximum of $C^1$ functions is locally Lipschitz
and has the stated right derivative, including ties. \\
62--70 & \pass & Positive gaps exclude inactive indices locally. Finiteness
makes the Taylor remainder uniform over the active set. Switching causes
no omitted case. \\
71--85 & \pass & Smoothness on $\Dom$, local existence, uniqueness and
continuous dependence require only finite real coefficients. No global
claim is inferred for all of $\Dom$. \\
86--96 & \pass & Face value is $c x_{i-1}$. The compact rectangle containing
$x,x^+$ and their segments lies within the open domain before any maximal time. \\
97--109 & \pass & Both $V>0$ and $V=0$ cases give $D^+V\leq LV$.
Lipschitz regularity justifies the almost-everywhere integrating-factor argument. \\
110--119 & \pass & Positivity claims have the required distinction between
$c>0$, $c=0$, and the zero initial state. Persistence is limited to the
maximal interval at this stage. \\
120--133 & \pass & The $b_i,d_i$ decomposition and both exponent limits in
variation of constants are exact. Coefficients are continuous on compact
time intervals, even though solution-dependent. \\
134--141 & \pass & Induction propagates the property ``positive for every
$t>0$,'' with no cumulative waiting times. Step-one connectivity suffices;
$q$ is irrelevant. \\
142--156 & \pass & $H=s b$ and the clipped plus-root formula are correct.
The radius is explicitly a comparison radius, not claimed optimal. \\
157--169 & \pass & The maximum inequality retains the correct $g$ loss;
the positive-root quadratic has constant term $-(a+c)$. \\
170--177 & \pass & The zero-radius inequality and the $s\geq R$ factor
bound both hold with the displayed signs. \\
178--189 & \pass & All three componentwise inequalities hold at every
active maximum. The negative $g$ multiplier reverses the reciprocal order. \\
190--195 & \pass & $b'<0$, the endpoint signs, quadratic multiplication,
and discriminant $(a+k)^2+4kc$ are correct. \\
196--203 & \pass & $k-a>0$ and the discriminant comparison correctly
select $R=0$ when $a+c\leq0$. Subtracting $b(R)=0$ is used only for $R>0$. \\
204--218 & \pass & Scalar domination yields the claimed global bound and
invariance of every cube of radius at least $R$. \\
219--228 & \pass & Exponential distance, rational zero-radius decay, and
both bounds independent of the initial state are correct. \\
229--240 & \pass & Whole-orthant attraction uses the second bounds, not
the initial-state-dependent exponential alone. Both outer entry times solve
the appropriate inequality exactly. \\
241--247 & \pass & The scalar interval barrier and compact continuation
are valid. The proof does not assume monotonicity of $H$. \\
248--261 & \pass & Positive-part comparison handles $W=0$ with or without
active coordinates. $H(X)-H(u)\leq LW$ is used only where $X>u$. \\
262--270 & \pass & Compact containment excludes the singular boundary;
bounded $F$ gives an endpoint limit and local continuation. No circular
use of already global vector solutions occurs. \\
271--282 & \pass & $w'=Ru'/u^2\leq-kRw$ and $0<w(0)<1$ give the
initial-state-independent bound. Initial values below $R$ are covered by
invariance. \\
283--292 & \pass & Reciprocal integration covers positive scalar data;
zero data use uniqueness. The distance bound and substitution of $T_r$
justify the stated quantifiers. \\
293--314 & \pass & The exact absorption equivalence is restricted to
$R>0$; the preceding nonentry assertion treats $R=0$ separately. \\
315--325 & \pass & $T_*=(\log2)/(kR)+(1+R)/g$ and the bounded-data time
have correct factors and stated positive-denominator hypotheses. \\
326--333 & \pass & Persistence rules out finite-time entry into $K_0$.
The coming-down estimate supplies $K_{2R}$ at $t_0$ uniformly. \\
334--343 & \pass & $X/(1+X)$ increases, so the negative strict-gap bound
is correct. Integrating until first entry and using invariance includes
the endpoint of the stated time. \\
344--347 & \pass & For the zero residue the inhibitory coordinate equals
$X$ at every maximum; the remaining time is $1/(hR)$. \\
348--355 & \pass & $\beta=(-a)_++(k+h)M+g$ and $\eta=c/(1+M)$ have
all required terms. The assumption $c>0$ makes $\eta>0$. \\
356--362 & \pass & Lower bounds on each loss and input yield the displayed
linear differential inequality. The exponential transformation has the
correct sign and base time. \\
363--371 & \pass & Repeated time-ordered integration gives $\tau^j/j!$.
There are $n$ retained terms, with largest power $n-1$, and a nonnegative
remainder. Wrapping indices does not invalidate it. \\
372--380 & \pass & The maximum at $t-\tau$ seeds one term for every target.
$d$ is positive and independent of initial data conditional on nonentry;
the text states that condition explicitly. \\
381--385 & \pass & The gap is $hRd$. The excess is at most $R$, so the
remaining time is $1/(hd)$, with no missing factor $R$ or $n$. \\
386--392 & \pass & Whole-orthant absorption implies the other forms.
For $g=h=0$, equal-coordinate data give exactly $s'=H(s)$ and cannot
reach $R$ at finite time. \\
393--398 & \pass & In the other failure case, $c=0$ preserves all zero
coordinates and $q\not\equiv0$ removes self-inhibition of the sole positive
one. Its logistic solution stays above $R$. The failure cases are exhaustive. \\
399--421 & \pass & The canonical corollary follows from the previously
verified hypotheses, with $R,c,g>0$. Neither coprimality nor an $h$-based
argument is needed for its absorption time. \\
422--427 & \pass & The archived expression for $B$ selects the correct
larger quadratic root. \\
428 & \minor & E01: ``the nonnegative root'' is ambiguous/incorrect as
a uniqueness description for $c=0,a>0$. Replace with ``the larger real
root, which is nonnegative.'' Formula unchanged. \\
429--434 & \pass & $H(B)=-c/(1+B)$ for $c>0$, so $B>R$; for $c=0$,
$B=R$. The claimed absorption consequences follow in both cases. \\
435--445 & \pass & Set invariance, absorption, attraction and stability
are distinguished correctly. Stability is relative to $\Orth$ and the
asymptotic-stability convention is expressly stated. \\
446--452 & \pass & The $k\geq0$ proposition asserts global existence and
the already qualified positivity conclusions, not the positive-$k$ radius. \\
453--466 & \pass & The exponential finite-time bound excludes blow-up.
The $a>0,k=c=g=h=0$ example correctly refutes a general absorbing-cube
guarantee for $k=0$. \\
467--475 & \pass & The $c<0$ face counterexample works for $n\geq2$.
The source explicitly excludes the $n=1$ case from this example. \\
476--479 & \pass & Ties, the loss sign, and coincident offsets are
accurately summarized. \\
480--484 & \pass & The zero radius need not solve the quadratic;
all divisions by $R$ are confined to positive-radius cases. \\
485--489 & \pass & The text correctly refuses to infer finite entry
merely from $H(R)=0$ and distinguishes the zero-radius obstruction. \\
490--495 & \pass & The dimension and connectivity relaxations are valid.
Instantaneous strict positivity is explicitly separated from uniform
long-time persistence. \\
496--504 & \pass & The strict coefficient assumptions are stronger than
needed for several conclusions. The negative-$k$ and negative-$h$
blow-up examples are exact existential counterexamples. \\
505--512 & \pass & Continuation uses compactness inside the smooth
domain. No later theorem or numerical experiment is a proof dependency.
The source remains DRAFT as requested. \\
\bottomrule
\end{longtable}
\normalsize

\section{Adversarial tests and reproducible arithmetic checks}
\subsection{Analytical attempts to break the claims}
\begin{enumerate}
\item \textbf{Very large data.} In the nonabsorbing scalar subcase
$c=g=h=0$, $a>0$, a synchronous state with $s_0>R=a/k$ satisfies
\[
 s(t)=\frac{R}{1-(1-R/s_0)e^{-kRt}}.
\]
Letting $s_0\to\infty$ recovers $R/(1-e^{-kRt})$ exactly as a
supremum. This verifies that the coming-down bound and whole-orthant
attraction are compatible with nonentry into $K_R$ at every finite
time. There is no hidden upper bound on admissible initial states.
\item \textbf{One supported coordinate.} With $c>0$, the integrating-factor
induction reaches every coordinate at any prescribed positive time,
even when the step-$q$ graph is disconnected. With $c=0,g=0,h>0$
and nonzero residue, exactly the same initial pattern supplies the
necessity counterexample. This stress test distinguishes the two regimes.
\item \textbf{Equal-coordinate initial states.} Direct substitution gives
$s'=H(s)-g s/(1+s)-h s^2$ for every $n,q$. At $s=R>0$, either
$g>0$ or $h>0$ makes the field strictly negative. At $g=h=0$ the
field vanishes there and uniqueness prevents finite-time crossing from
above. This substitution is only a test family, not a P02 equilibrium audit.
\item \textbf{Ties and switching.} For $n=5,q=2,a=k=c=g=h=1$ and
$x=(2,2,0,0,0)$, the active derivatives are $F_0=-4$ and
$F_1=-10/3$. Thus $D^+X=-10/3$, not the derivative of an arbitrarily
selected active coordinate. The finite-maximum expansion proves the
same rule at an actual switch, regardless of how often the active set
changes. The elementary functions $y_1(t)=t,y_2(t)=-t$ also test the
right-derivative assertion at a switch: it is $1$ at zero although
the two-sided derivative does not exist.
\item \textbf{Requested dimensions and offsets.} The scalar reduction
explicitly tests $n=1$; the coincident-neighbor case tests $n=2,q=1$.
For $n=5$, every residue $0,1,2,3,4$ satisfies the same maximum and
propagation estimates. For even $n=6,10$, the cases $q=1,n-1,n/2,0$
are covered with all coincident terms retained; $n=6,q=2$ additionally
tests absence of a coprimality requirement. The exact calculation suite
below checks every residue for all five dimensions.
\item \textbf{Parameter edges.} $c=0$, $g=0$, $h=0$, $k=0$,
$a+c=0$, $a+c<0$, and small positive $R$ have been treated explicitly
above. The strongest alleged loophole, loss of a uniform positive
coordinate floor near the zero state, is avoided by conditioning that
floor on continued nonentry. No unconditional persistence claim is proved.
\end{enumerate}

\subsection{Exact rational checks: evidence and limits}
A deterministic check using exact rational arithmetic was run independently
from the defining right-hand side. It performed no time integration,
used no floating-point tolerances, and is not offered as a proof over
an infinite state or parameter space. All asserted inequalities passed:
\begin{center}
\begin{tabular}{lr}
\toprule
Check & Count \\
\midrule
State/parameter/offset cases for the maximum inequality & 7,956 \\
Active-coordinate inequalities within those cases & 19,682 \\
Constructed positive-root identities & 48 \\
Coordinate inequalities for the propagation lower bound & 82,080 \\
Failed assertions & 0 \\
\bottomrule
\end{tabular}
\end{center}
The following complete recipe permits reproduction without relying on
an undistributed numerical data file. Let $\epsilon=10^{-12}$ and
$L=10^{12}$. Use $n\in\{1,2,5,6,10\}$ and integer representatives
$q=0,1,\ldots,n$, with $q=n$ deliberately checking the zero-residue alias.
For each $n$, use these $8+2n$ state entries, retaining repetitions:
\begin{itemize}
\item all-equal states at $0,\epsilon,1,2,L$;
\item for each coordinate, the two states supported there with value $1$ or $L$;
\item $x_i=2$ for even $i$, zero for odd $i$; $x_i=(i+1)/n$;
and $x_i=2$ for $i<2$, zero otherwise, using indices $0\leq i<n$.
\end{itemize}
Take the Cartesian product with the following 13 coefficient tuples:
\begin{center}
\begin{tabular}{ccccc@{\qquad}ccccc}
\toprule
$a$&$k$&$c$&$g$&$h$&$a$&$k$&$c$&$g$&$h$\\
\midrule
$1$&$1$&$1$&$1$&$1$&$-2$&$1$&$1$&$1$&$1$\\
$1$&$1$&$0$&$1$&$0$&$-1$&$1$&$1$&$0$&$0$\\
$1$&$1$&$1$&$0$&$1$&$0$&$1$&$\epsilon$&$\epsilon$&$\epsilon$\\
$1$&$1$&$0$&$0$&$1$&$\epsilon$&$1$&$0$&$\epsilon$&$0$\\
$1$&$1$&$1$&$0$&$0$&$1$&$\epsilon$&$1$&$1$&$1$\\
$1$&$1$&$0$&$0$&$0$&$-1$&$1$&$2$&$0$&$1$\\
$0$&$0$&$1$&$1$&$1$&&&&&\\
\bottomrule
\end{tabular}
\end{center}
At every active index, evaluate \eqref{A01:field} directly and assert
$F_i\leq H(X)-gX/(1+X)\leq H(X)$. The $k=0$ tuple tests this
inequality only, not a radius formula requiring $k>0$.

For the 48 root cases, take
\[
 R\in\{\epsilon,1/10,1,L\},\quad
 k\in\{\epsilon,1,10\},\quad
 c\in\{0,\epsilon,1,10\},\quad
 a=kR-\frac{c}{1+R}.
\]
Check $a+c>0$, the quadratic residual is zero, and
\[
 (a+k)^2+4kc=(2kR-a+k)^2,\qquad 2kR-a+k>0.
\]
This verifies selection of the plus root without approximating a square
root. Check \eqref{A01:factor} exactly at $s=R,2R,100R$.
For the lower-bound checks, retain these root cases, the same dimensions
and representatives $q$, put $M=2R$, use $g,h\in\{0,\epsilon,1\}$,
and take $x_i=M(i+1)/n$. At every coordinate assert
$F_i\geq-\beta x_i+\eta x_{i-1}$ with \eqref{A01:constants}.
The analytical proofs, rather than this finite suite, establish the
factorial propagation estimate, finite-time classification, and all
limiting and uniform quantifiers.

\section{Hypotheses, corrections, and downstream consequences}
\subsection{Unnecessary hypotheses for particular conclusions}
\begin{longtable}{@{}P{0.30\textwidth}P{0.65\textwidth}@{}}
\toprule
Canonical restriction & What the audited proof actually needs \\
\midrule\endhead
$n\geq5$ & A fixed positive integer $n$ suffices for every P01 theorem.
The negative-$c$ boundary counterexample specifically needs $n\geq2$. \\
$1\leq q\leq n-1$ & Any residue works; zero is handled explicitly in
the extended absorption criterion. \\
Coprimality or distinct offsets & Neither is needed for any P01 result.
Coprimality was not imposed by P01 and must not be silently added. \\
$a>0$ & $a\in\RR$ suffices for the extended results. $a+c>0$ is exactly
the condition $R>0$ under the other sign assumptions, not a condition
for global existence. \\
$k>0$ & Not needed for local well-posedness, nonnegative invariance,
or finite-time positivity. $k\geq0$ suffices for global existence in
the stated nonnegative-loss class. The particular radius and coming-down
bounds require $k>0$. \\
$c>0$ & $c\geq0$ suffices for invariance and the given comparison bounds.
For $n\geq2$, $c>0$ is necessary to guarantee immediate strict positivity
from every nonzero nonnegative initial state; for $n=1$, it is unnecessary
for that implication. \\
$g>0$ & $g\geq0$ suffices for the stated radius and bounded-attraction
results. Strict positivity is only one sufficient absorption mechanism;
the criterion identifies the alternatives. \\
$h>0$ & $h\geq0$ suffices for the radius and bounded-attraction results.
Strict $h>0$ is unnecessary for canonical finite entry when $g>0$. \\
Signs of $a,k,g,h$ in the local face and positivity arguments & None
is needed there: finite real coefficients and $c\geq0$ suffice up to
the maximal time. Subsequent global bounds use their own hypotheses. \\
\bottomrule
\end{longtable}
This is a statement-by-statement inventory of unused canonical
restrictions, not an optimal classification for arbitrary signed
coefficients. In particular, the counterexamples outside the proved
domain do not turn sufficient sign assumptions into necessary ones
for every individual equation.

\subsection{Hidden-assumption check and suggested clarifications}
No missing mathematical assumption was found. The following operational
qualifications are explicit or inherent in the source and must travel
with any later quotation of its results:
\begin{enumerate}
\item $n$ is a positive integer, coefficients and states are finite,
the coefficients are homogeneous and time independent, and initial
states are nonnegative for global forward claims.
\item The domain $\Dom$ supplies local solutions; the global bounded
flow is established on $\Orth$. Smoothness is not an assertion that
arbitrary negative initial data stay in $\Dom$ forever.
\item Uniform entry times are uniform over initial states for fixed
parameters, $n$, and $q$. In the propagation proof, choose one fixed
$\tau>0$ independent of initial data.
\item The floor $d$ applies only while entry has not occurred. It neither
covers the zero trajectory nor establishes uniform persistence after entry.
\item Positive invariance supplies all times after first entry. Scalar
local Lipschitz regularity and absolute continuity justify integration
of the Dini inequalities; differentiability of $X$ everywhere is not assumed.
\item Set stability and attraction are relative to $\Orth$ and use
distance to the cube, under the expressly stated convention for
asymptotic stability. None is an equilibrium stability statement.
\end{enumerate}
Apart from required correction E01, the following are optional clarity
improvements, not new gaps: spell out $n\in\ZZ_{\geq1}$ in (P01-W);
append ``for fixed parameters and dimension'' to the whole-orthant
uniformity claim; and call \eqref{A01:d} a ``conditional pre-entry
floor'' when quoting it elsewhere. Do not replace a proved theorem
with a weaker assertion merely because its strong uniform quantifier
is unusual. Do not enlarge $R$ back to $B$ as a purported repair: the
sharper radius is already correct.

\subsection{Downstream consequences without proceeding to P02}
\begin{enumerate}
\item \textbf{Available foundation.} Later demand-modulated modules may
use the canonical smooth vector field, unique global nonnegative
forward solutions, immediate positivity from nonzero data, and the
invariant absorbing cube with the stated radius and time bound. E01
does not invalidate any of these dependencies.
\item \textbf{What remains unproved downstream.} P01 does not classify
synchronous equilibria, prove existence of a positive equilibrium,
establish equilibrium stability, compute a spectrum, or prove a
bifurcation. Those are separate proof obligations. No such work was
undertaken in this audit.
\item \textbf{Using the box correctly.} Every nonnegative equilibrium
of this field must lie in $K_R$, since a stationary state outside it
would contradict attraction. This is a permissible necessary bound,
not an equilibrium existence or uniqueness proof. Set attraction does
not imply convergence of a trajectory to an individual equilibrium;
nontrivial dynamics inside the box are not excluded.
\item \textbf{Parameter boundaries.} If a later argument sets $c=0$,
it loses general positivity propagation. If it sets $g=0$, it must
use the remaining absorption criterion, not the formula containing
$1/g$. If $a+c\leq0$, $K_R=\{0\}$ is attracting but not finitely
entered from nonzero data. If $k=0$, the radius formula and the present
bounded-attraction conclusions are unavailable without a new argument.
\item \textbf{No transfer to the matched bilinear comparator.} The
positive input $c x_{i-1}/(1+x_i)$ is central to both the bound and
boundary propagation. P01's radius, positivity propagation, and global
boundedness cannot be copied to the different bilinear equation merely
because the coefficient names match. Any such use in P08 needs its
own hypotheses and proof.
\item \textbf{No uniform persistence or optimality claim.} The lower
floor in Theorem 5.6 is temporary and conditional. Neither a permanent
positive coordinate bound nor a smallest invariant cube has been established.
Likewise no heterogeneous, delayed, stochastic, or time-varying version
has been certified by this audit.
\end{enumerate}

\textbf{Disposition.} The canonical foundation is verified for the exact
source hash identified above. The extended theorems and their complete
absorption classification are mathematically verified, but the exact
extended text receives MINOR REVISION solely for E01. The audit changes
no P01 source, status marker, registry, compendium, or later module.
The delivered archive contains exactly one member,
\path{audits/A01_wellposedness_and_absorbing_box.tex}, a standalone
LaTeX document. Compilation and archive checks validate the deliverable;
the mathematical verdicts are justified by the independent arguments
in this document.
\end{document}
